\chapter{Conclusion and Future Work}

\section{Project Summary}
This project successfully architected and implemented a high-performance, industrial-grade acoustic anomaly detection system. By transitioning from traditional spectral feature engineering to a sophisticated \textbf{Dual-Track} pipeline, we provided a solution that balances sub-10ms edge inference with deep cloud diagnostics.

Key achievements include:
\begin{itemize}
    \item \textbf{High-Fidelity Diagnostics:} Achieved a validated \textbf{95.61\% test accuracy} and \textbf{0.9267 Macro F1-score} using the \textbf{EfficientAT} architecture, maintaining high precision even under rigorous cluster-aware evaluation.
    \item \textbf{Edge Performance:} Developed a native Rust feature extraction engine (\texttt{fast\_features}) providing a \textbf{57x speedup} over standard libraries, enabling a \textbf{9.8 ms} end-to-end pipeline latency for edge-based XGBoost inference.
    \item \textbf{Generative Data Augmentation:} Addressed critical industrial class imbalance by synthesizing high-fidelity abnormal signatures via a \textbf{Conditional GAN (CGAN)}, boosting minority class F1-scores by over 12\%.
    \item \textbf{Generalized Robustness:} Eliminated recording session bias through a \textbf{DBSCAN Cluster-Aware} evaluation regime, ensuring the system generalizes to new factory environments.
\end{itemize}

\section{Key Takeaways}
The defining success of this project was the resolution of the "Generalization Trap." By treating recording session bias as a manifold clustering problem rather than a standard classification task, we moved beyond superficial high-accuracy metrics to a truly robust system. The deployment of a hybrid architecture—leveraging the speed of gradient-boosted trees for continuous monitoring and the depth of temporal-attention networks for diagnostic deep-dives—represents a state-of-the-art approach to industrial IoT monitoring.

\section{Future Work}
\begin{itemize}
    \item \textbf{Federated Learning:} Implementing a federated training framework to allow models to learn from across multiple factory sites without exposing sensitive acoustic data, improving global robustness while maintaining local privacy.
    \item \textbf{Unsupervised Outlier Exposure:} Integrating an unsupervised branch (e.g., Variational Autoencoders) to detect "zero-day" mechanical failures that were never seen during the training phase.
    \item \textbf{On-Device Quantization:} Utilizing INT8 quantization and ONNX Runtime to deploy the EfficientAT diagnostic model directly onto ARM-based edge gateways, reducing the dependency on cloud connectivity.
\end{itemize}
